\textbf{Fuente del ciclo laboral y habilitación.}
Synaptix selecciona un Registro Laboral y de Habilitación (RLH) nuevo, dentro de Auditoría y seguridad y del Runtime Operacional de Borde. El RLH registra la decisión de negocio; la Autoridad local de acceso ejecuta su efecto sobre identidades, sesiones y permisos. Ambos usan el almacenamiento y la réplica síncrona del clúster local. Los tres núcleos cloud consumen sus eventos por el proceso de sincronización y conservan la versión de acceso aplicada. Esta incorporación mantiene los nueve dominios y los cinco procesos especializados.

Gerencia conserva la competencia sobre el personal propio y Administración registra sus decisiones mediante titular y suplente delegados. Cada una de las 62 empresas contratistas registra representantes titulares y suplentes con facultad sobre sus vínculos; Varadero verifica empresa, representación y acreditación. Synaptix administra directorio, integración, claves, respaldo, alarmas y aplicación automática. El responsable funcional registra una decisión laboral; el efecto técnico no espera un ticket ni una segunda digitación. Los mandatos, las identidades nominativas y sus suplencias se incorporarán al registro de habilitación del servicio; no se presentan como acuerdos ya suscritos.

La Tabla~\ref{tab:rlh-cobertura} compara origen y autorización por estado de cobertura. Las responsabilidades y el flujo se desarrollan en el texto; la matriz permite localizar cuándo puede emitirse un permiso.

\begin{longtable}{L{\dimexpr.28\linewidth-2\tabcolsep\relax}L{\dimexpr.32\linewidth-2\tabcolsep\relax}L{\dimexpr.40\linewidth-2\tabcolsep\relax}}
\caption{Origen, cobertura y aplicación del registro laboral}\label{tab:rlh-cobertura}\\
\hline
\EncabezadoTabla{Situación} & \EncabezadoTabla{Origen competente} & \EncabezadoTabla{Nueva emisión}\\ \hline
\endfirsthead
\hline
\EncabezadoTabla{Situación} & \EncabezadoTabla{Origen competente} & \EncabezadoTabla{Nueva emisión}\\ \hline
\endhead
Propio, cobertura vigente & Gerencia / Administración delegada & Permitida tras autenticación y comprobación de vínculo\\ \hline
Contratista, cobertura vigente & Representante con mandato & Permitida dentro del alcance registrado\\ \hline
WAN perdida, LAN sana & Responsable competente local & Permitida solo con circuito de origen vigente\\ \hline
Baja exterior desconocida & Origen exterior no disponible & Denegada por pérdida de cobertura\\ \hline
Terminal incomunicado & Ningún origen alcanzable & Denegada; conserva solo permiso residual\\ \hline
\end{longtable}
\FuenteTabla{Diseño de Synaptix; \parencites[art.~22, p.~16]{pucv-admin}[RT-12.10/13, p.~26]{pucv-transversal}[cap.~15, RT-03.10/13]{caso07}.}

La disponibilidad del RLH por LAN permite emisión durante una caída exterior solo cuando permanece vigente el circuito de su origen competente. En las dos últimas filas, conservar el registro protege evidencia, pero no aporta esa autorización nueva. Cada decisión conserva persona, empleador, vínculo, representante y mandato, identificador de evento, secuencia por empleador, versión del vínculo, tipo, alcance, evidencia y firma. Se separan decisión, efecto laboral real, recepción y aplicación. Una transacción durable registra decisión, estado y evento de salida; baja o suspensión establece denegación antes del acuse, incrementa la versión y revoca las familias de sesión y refresco. Alta o cambio provisiona automáticamente el alcance aprobado. Una reactivación necesita una decisión competente posterior vinculada al estado anterior; un mensaje antiguo nunca vence una baja.

El objetivo propio de aceptación es aplicar y propagar el cambio en componentes alcanzables en un máximo de 60 segundos desde recepción, con bloqueo en el RLH antes del acuse. El límite de baja se mide desde la desvinculación real. Para decisiones externas, \(D\) es la demora desde efecto hasta recepción, \(P\) la aplicación, \(L\) la vida residual máxima admitida y \(\varepsilon\) la incertidumbre temporal: se exige \(D+P+L+\varepsilon\leq24\,h\). Un permiso nuevo de 24 horas añadido a una entrega tardía no satisface esa cota. Durante la pérdida de cobertura no se renueva desde una copia antigua. Continúa pendiente desarrollar la cobertura completa de decisiones exteriores bajo aislamiento total; esta selección de fuente no elimina esa brecha.

\textbf{Sesión común y credenciales.}
Auditoría y seguridad incorpora un registro común de sesiones con identificador de sesión, identidad, dispositivo, emisor, autenticación, expiración absoluta, última actividad, versión de acceso y estado. En el recinto se almacena y replica con el RLH; en nube se mantiene en el esquema de seguridad y se aplica por un módulo común en los tres núcleos. Cognito verifica la identidad federada con enlace disponible. Tras esa autenticación se crea una sesión propia de Synaptix, con identificador nuevo y alcance del RLH, para todos los módulos autorizados. La Autoridad local verifica el autenticador individual inscrito durante una caída exterior.

La sesión interactiva dura como máximo ocho horas desde autenticación. Caduca por quince minutos de inactividad en oficina o portal y diez minutos en móvil. El bloqueo conserva la cola cifrada y exige autenticación para volver a interactuar. Una nueva autenticación reemplaza las credenciales y el identificador anterior. Se admite una sesión operacional por persona y dispositivo y una sola identidad activa en un terminal compartido; el cambio de usuario revoca la sesión anterior y limpia su vista. Una sesión técnica privilegiada se registra por separado y no concede facultades operacionales.

La credencial de sesión se firma y dura como máximo cinco minutos; conserva emisor, audiencia, sujeto, dispositivo, alcance, versión, identificador, emisión y vencimiento. El refresco es aleatorio, de un solo uso, rotatorio y vinculado a la sesión y al dispositivo. El registro conserva su huella y el estado consumido para detectar reutilización; esta revoca toda la familia y genera una incidencia. La rotación no extiende las ocho horas ni la vigencia del vínculo o permiso. Las claves de firma se protegen y rotan con versiones verificables; la clave local no se presenta como una clave de Cognito.

En un terminal completamente desconectado, la autenticación individual abre una sesión local dentro del permiso operacional aún válido. El verificador local aplica las mismas caducidades y genera credenciales breves firmadas por la clave técnica inscrita del dispositivo; el refresco local se registra como consumido en una transacción durable. Su audiencia es exclusivamente el verificador local: nube y Runtime rechazan esa credencial como sesión propia o prueba suficiente para emitir un nuevo permiso. La sincronización utiliza la identidad técnica del dispositivo y la evidencia de autorización vigente en la actuación capturada. Ningún refresco local renueva el permiso operacional o el vínculo laboral. Se rechazan reloj no fiable, desafío repetido, baja conocida y permiso vencido. El límite de 24 horas corresponde al permiso operacional y no a la credencial de sesión.

Los identificadores y tokens se transportan por cabecera de autorización o cookie Secure/HttpOnly/SameSite según cliente, bajo TLS y protección CSRF para cookies. No se incorporan a rutas, consultas URL, nombres de archivos, registros técnicos ni Referer. El cierre de sesión o la baja revoca el identificador y la familia en el registro común antes del acuse. Todos los módulos y emisores alcanzables comprueban ese estado antes de confirmar acciones; la caché de autorización del gateway se deshabilita para estos accesos. Un dispositivo aislado aplica la revocación al recibirla y su límite firmado mientras carezca de canal. La propagación inmediata a un dispositivo físicamente incomunicado sigue siendo una limitación de cobertura de RT-12.02/07.

\textbf{Acceso de emergencia, custodia y control.}
Se seleccionan dos identidades individuales de último recurso, asignadas al responsable de Seguridad del servicio y al responsable técnico de Operación como titular y suplente. Cada identidad se enrola por separado con MFA y verificación local de la clave pública. La nominación se registra antes de habilitación; no existe una contraseña compartida que permita atribuir al portador la identidad del operador.

La credencial de recuperación y el autenticador se custodian por separado en depósitos cifrados y sobres precintados. La apertura requiere dos custodios distintos: Seguridad de Synaptix y la contraparte funcional autorizada por Gerencia. El registro conserva custodio, operador, incidente, motivo, aprobaciones, instante, identificador y huella del sobre, entrega y devolución. La doble custodia autoriza la apertura; el operador debe además acreditar su identidad y habilitación laboral vigente contra el RLH competente. Un sobre o una copia antigua no sustituyen esa comprobación.

La activación requiere un incidente crítico registrado y dos autorizaciones por roles diferentes. La sesión dura como máximo 60 minutos y permite restaurar el servicio, recuperar el registro y revocar credenciales. Su política impide altas laborales, ampliación de facultades de negocio, extracción masiva y borrado de auditoría. Se conservan localmente el comienzo, cada comando y efecto, resultado y cierre, con evidencias protegidas y envío posterior. Si no puede mantenerse la auditoría o acreditarse la habilitación del operador, se deniega la sesión.

Si el RLH no está disponible, la recuperación inicial usa la consola física del nodo y un medio técnico de arranque protegido y firmado, preparado con la versión Ubuntu Server de la BOM y la tarea de restauración de Synaptix. Su autorización de apertura y sus claves se custodian fuera del RLH, con dos custodios e identificación individual del ejecutor. El medio verifica su firma y el manifiesto de la copia, ejecuta únicamente restauración e integridad y registra localmente la cronología en un depósito separado. No obtiene una sesión del registro que intenta recuperar. Restaura bajo esa custodia una copia cifrada y autenticada del RLH y sus claves, verifica el punto de recuperación y mantiene las aplicaciones bloqueadas hasta comprobar la vigencia del origen y del operador. La recuperación física de una copia no concede una sesión de aplicación ni confirma por sí sola actualidad laboral.

Al cerrar se revocan sesión y familia, se rotan todas las credenciales utilizadas, se reprecintan los depósitos y Seguridad encarga una revisión independiente con informe de incidente y acciones correctivas. La aceptación previa y el ejercicio semestral incluyen ausencia del titular, caída cloud, baja del operador, negativa de un custodio, clave inválida, salida del alcance, falla de auditoría y cierre/rotación. Se conserva acta con identidades, cronología, resultados y correcciones. Este procedimiento desarrolla RT-12.13; no se utiliza para extender permisos ni acreditar las autonomías de operación \parencite[RT-12.02/07/08, p.~25; RT-12.10/13, p.~26]{pucv-transversal}.
